\section{Claim types and the minimal formal setting}
\label{sec:claims}
% ============================================================================

This section fixes the two pieces of infrastructure that everything else
depends on: a typing discipline for claims, and the small stock of standard
finite mathematics in which all results are stated.

\subsection{Four claim lanes}
\label{sec:lanes}

Every scientific sentence in this paper belongs to exactly one lane.

\begin{description}
\item[Tier S (structural).] Mathematical content about declared finite
objects. Tier~S subdivides by \emph{claim form}: a \emph{theorem-grade exact
result} is proved or exhaustively computed on its declared carrier; a
\emph{schema} is a definition or conditional statement whose hypotheses are
displayed and not discharged; a \emph{pointer} re-scopes an existing result to
a new reading without new mathematics; a \emph{composition} combines landed
results under a declared convention. Tier~S claims are established relative to
their carriers and interfaces, and to nothing else.
\item[Tier R (recognition).] A proposed mapping between a structural object
and a familiar mental or phenomenal category, for example reading a
constraint's dual price as felt valence. Recognition bridges are hypotheses.
They are \emph{proposed}, \emph{named}, and \emph{retractable}; they are never
derivations, and a counterexample to a bridge retracts the bridge, not the
mathematics under it.
\item[Tier O (ontological).] A claim about what satisfaction of the structure
would be. All Tier~O content in this paper is conditional on named postulates
(Section~\ref{sec:t17}); the paper contains no unconditional Tier~O sentence.
\item[OPEN.] A target that is stated, motivated, and not achieved. Open
targets are recorded in the artifact with explicit markers rather than
implied by silence.
\end{description}

Table~\ref{tbl:grades} summarizes the lanes, their evidential support, and the
inference each lane blocks. Figure~\ref{fig:lanes} shows the same discipline as
a flow constraint: computational and literature inputs enter only their proper
lanes, and no arrow runs from computation or literature directly to
realization.

\begin{table}[t]
\centering\small
\begin{tabular}{p{0.1\linewidth}p{0.2\linewidth}p{0.22\linewidth}p{0.18\linewidth}p{0.2\linewidth}}
\toprule
Grade & Object & Evidential support & Licensed wording & Blocked inference \\
\midrule
\gradeS{} theorem/ exact & declared finite carrier and result & proof, or exhaustive exact enumeration with tests; selected Lean kernels & ``proves on the declared carrier,'' ``computes exactly,'' ``exhaustively enumerates'' & generalization beyond the carrier; any empirical or phenomenal conclusion \\
\gradeS{} schema & conditional statement with displayed hypotheses & definition-level checking; conditional proofs & ``establishes under assumptions,'' ``gives a conditional schema'' & unconditional use; silent discharge of hypotheses \\
\gradeR{} & proposed structure-to-phenomenon mapping & argued fit of profile; no derivation & ``proposes,'' ``recognizes,'' ``interprets, conditional on'' & treating the mapping as derived or measured \\
\gradeO{} & conditional ontological implication & named postulates plus a mechanized conditional theorem & ``implies after assuming,'' ``is conditional on,'' ``purchases'' & detaching the conclusion from its premises \\
\gradeOpen{} & stated, unachieved target & explicit negative markers in the artifact & ``remains open,'' ``is not exhibited'' & ``follows in principle,'' ``is nearly proved'' \\
\bottomrule
\end{tabular}
\caption{\textbf{Tbl-1: claim-grade legend.} Each grade names its object, its
support, its licensed verbs, and the inference it blocks. Scope/price: the
legend is a reporting discipline for this paper, not a proposed norm for the
field.}
\label{tbl:grades}
\end{table}

\begin{figure}[t]
\centering
\fbox{\parbox{0.92\linewidth}{\small\emph{[Figure F1: claim and evidence
lanes; rendered per the artifact figure plan.]} Four horizontal lanes:
structural result, recognition bridge, ontological implication, open
obligation. Python enumeration and Lean declarations feed only the structural
lane; philosophical postulates feed only the ontological lane; literature
feeds comparison, never evidence. No arrow connects computation or literature
directly to realization.}}
\caption{\textbf{F1: claim and evidence lanes.} Takeaway: every sentence has
one lane, and evidence types cannot cross lanes. Scope restriction: the
diagram is schematic; the sentence-level assignment lives in the artifact's
claim ledger.}
\label{fig:lanes}
\end{figure}

\subsection{The minimal formal setting}
\label{sec:setting}

All objects are finite. Probabilities and costs are exact rationals; no
floating-point value appears in any result.

\begin{definition}[Carrier, actions, lens]
A \emph{carrier} is a finite set $X$ of states. A \emph{transition/action
family} over $X$ is a finite set $A$ of actions with, for each $a\in A$,
either a map $F_a : X \to X$ or a row-stochastic kernel $P_a$ on $X$; whether
dynamics are deterministic or stochastic is declared per construction. A
\emph{lens} (observation map) is a map $q : X \to Y$ onto a finite value set;
its fibers form the \emph{observational partition} $\Pi_q$ of $X$. A
\emph{policy for $q$} is a feedback rule that selects actions measurably with
respect to $\Pi_q$ (more generally, with respect to the history of $q$-values).
\end{definition}

\begin{definition}[Quotient, refinement, factorization defect]
A partition $\Pi'$ \emph{refines} $\Pi$ when every block of $\Pi'$ is contained
in a block of $\Pi$. Given two lenses $\pi_0, \pi_1$ on the same carrier, the
\emph{factorization defect} is
\[
\Dfact(\pi_0,\pi_1) \;=\; \{(s,s') : \pi_0(s)=\pi_0(s') \text{ and }
\pi_1(s)\neq\pi_1(s')\},
\]
the set of \emph{split pairs}: pairs identified by the coarse description and
separated by the fine one. $\pi_1$ factors through $\pi_0$ (there is $\varphi$
with $\pi_1=\varphi\circ\pi_0$) if and only if $\Dfact(\pi_0,\pi_1)=\emptyset$.
A nonempty defect is the exact, checkable meaning of ``the fine content is not
recoverable from the coarse description'' used throughout this paper.
\end{definition}

\begin{definition}[Predictive and observational equivalence]
Fix a carrier with declared dynamics and a lens $q$. Two states (or histories)
are \emph{observationally equivalent} when all admissible current observables
agree on them; the induced quotient is the \emph{current public quotient} $Q$.
They are \emph{predictively equivalent} when they agree on the full family of
admissible \emph{future} behavior (continuation laws, or distributions over
future observable trajectories under all admissible continuations); the induced
quotient is the \emph{predictive quotient} $M$. Predictive equivalence refines
observational equivalence, giving a canonical comparison map $\pi : M \to Q$.
The gap between them, same $Q$-class with different $M$-class, is a
factorization defect and is the recurring engine of this paper.
\end{definition}

Predictive quotients are standard mathematics: they are the finite controlled
analog of causal-state equivalence in computational mechanics
\citep{CrutchfieldYoung1989, ShaliziCrutchfield2001}, and we import the
construction, not the stationary-process theorems (the distinction matters in
Section~\ref{sec:t1}).

\begin{definition}[Closure operator and audit]
A \emph{closure operator} on a finite structured set is an extensive, monotone,
idempotent endomap; its fixed points are the descriptions the structure
recognizes as complete. When a closure-like operator $E$ is only approximately
idempotent, its \emph{idempotence defect} $\delta(E)=\tfrac12\max_z
\lVert E^2(z,\cdot)-E(z,\cdot)\rVert_1$ is computed exactly and reported rather
than assumed small. An \emph{audit} is a declared, executable check on a
construction, such as a defect computation, a control run, or an exhaustive
comparison, whose outcome is recorded with the result. A \emph{gate} is a named audit
whose passage is required before a claim may be graded; the gate vocabulary
used here (sufficiency, descent, stability, control-of-content, no-smuggling,
visibility, audit, strictness) is imported from the interaction calculus of
the companion corpus (Appendix~\ref{app:imports}) and is used as a reporting
format, not re-proved.
\end{definition}

\begin{definition}[Viability kernel]
Fix a carrier $X$, a safety predicate, feasible action sets, and declared
successor supports. The \emph{viability kernel} $\Kviab$ is the greatest fixed
point of the controlled-invariance operator
\[
V(K) = \{\, s \in K : s \text{ safe and } \exists a \in A_{\mathrm{feas}}(s)
\text{ with } \mathrm{Post}(s,a) \subseteq K \,\},
\]
computed by descending iteration; this is the standard discrete-time viability
kernel \citep{AnayaEtAl2009Viability, AubinBayenSaintPierre2011}. For a lens
$\ell$, we write $\Kviab(\ell)$ for the best kernel achievable by
$\ell$-measurable policies. A nonempty kernel means controlled persistence in
the model, nothing more.
\end{definition}

\subsubsection*{The trust chain}

Every result travels the same four-step chain, and the paper reports which
steps each result completed. (1)~A \emph{mathematical statement} on a declared
carrier, with grade and nonclaims, in the theorem catalog. (2)~A \emph{Python
enumeration}: exact rational computation of every quantity, with declared
controls, under test. (3)~A \emph{Lean declaration}: a named kernel per law,
whose scope boundary (what is proved versus what is pinned as certified
constants matching the Python computation) is part of the result.
(4)~An \emph{interpretation}, which is where lanes R and O begin and where the
ledger's blocked-inference column applies. Prose never outranks an executable
artifact: the theorem catalog controls statements, the current Python
computations control values, the Lean sources control what is machine-checked.

\subsection{Six Birds vocabulary, second}
\label{sec:sixbirds}

The companion corpus organizes the same mathematics under a six-role grammar
(threshold, budget, route, staging, packaging, and audit; P1--P6) with a typed
calculus of promotion between description layers: a \emph{promotion bridge}
carries content from a lower to a higher layer; \emph{strict} promotion means
the higher description does not factor through the lower one
($\Dfact\neq\emptyset$); \emph{no-smuggling} and \emph{no-overreading} rules
police what may be used and claimed across layers; and \emph{same-level
self-audit} is provably never accepted, while level-shifted, partial audit
is. This paper uses those results as imported theorems with exact citations
(Appendix~\ref{app:imports}); the six roles never organize sections, are never
renamed as neurocognitive faculties, and impose no requirement that every role
be active in every carrier. Analogy with a sibling result is never treated as
theorem transport.

% ============================================================================
